\section{Introduction}\label{sec:intro}

A process in time is often accessible only through \emph{windows}: overlapping
finite views, each recording the joint behaviour of finitely many coordinates over
finitely many steps. The resulting local probability tables may agree wherever they
overlap yet be generated by no single underlying process --- the observational face
of contextuality~\cite{AbramskyBrandenburger2011}. We ask, for symbolic dynamical
observation, when local coherence forces global realisability, and answer it for a
large structured class.

A finite alphabet $A$ and finite supports $F\subseteq\bbZ$ give \emph{contexts},
each a partition of $A^F$; a \emph{protocol} is a family of contexts, and
window-data assigns probabilities to the cells of each context.

\begin{definition}[$\EA$, $\PR(\Rc)$, commensurability]\label{def:ea-pr}
Window-data is \emph{empirically adequate} ($\EA$) if it is pairwise coherent:
the contexts agree on the common subalgebra of each pair. It is
\emph{probabilistically realisable} in a class $\Rc$ ($\PR(\Rc)$) if some member
of $\Rc$ (all measures / stationary / supported on a subshift / memory-$m$
Markov, an observer's commitment lattice) induces it. A \emph{model} is
\emph{contextual} when it is $\EA$ but not $\PR(\Rc)$. A protocol is
\emph{commensurable} when $\EA \Rightarrow \PR(\Rc)$ for all coherent data.
\end{definition}

The two protocol axes are \emph{reporting} (window layout) and \emph{coarsening}
(compression); the commitment class $\Rc$ is the third, graded, axis.

On a finite outcome set the coherent data form a polytope $C$ (the local/marginal
polytope~\cite{WainwrightJordan2008}) and the realisable data the polytope
$R=\mathrm{conv}\{v(x)\}$, the convex hull of the cell-incidence vectors of the
global configurations, with $R\subseteq C$.

\begin{proposition}[commensurability is $C=R$]\label{prop:cr}
A protocol is commensurable iff $C=R$. A contextual model is a point of
$C\setminus R$; over rational data any separating facet certifies
non-realisability.
\end{proposition}

Both polytopes are enumerable in exact arithmetic, so a commensurability verdict
is a proof that $C=R$ or an exhibited rational witness in $C\setminus R$
(\S\ref{sec:method}).

The polyhedral facts are classical: $C$ and $R$ are the marginal polytopes of
Abramsky--Brandenburger's presheaf model~\cite{AbramskyBrandenburger2011}; acyclic
extendability
is Vorob'ev~\cite{Vorobev1962}; the constrained case is
De~Loera--Onn~\cite{DeLoeraOnn2004} and
Pr\r{u}\v{s}a--Werner~\cite{PrusaWerner2015}; the parity mechanism is stable-set
integrality~\cite{NemhauserTrotter1974,Chvatal1975,HoffmanKruskal1956}. Over this
sits an observational and dynamical layer: a cyclic protocol read as recurrence,
so that a length is contextual exactly when the layered ring carries a
winding-$\ge2$ orbit (\S\ref{sub:winding}), and the tamings assembled into a
localisation of the polytope (\S\ref{sub:tamings}). The winding criterion and the
localisation assembly sit above the satisfiability-width line of
Barto--Kozik~\cite{BartoKozik2014} and the tightness criterion of
Thapper--\v{Z}ivn\'y~\cite{ThapperZivny2016}; the signable tamings are the
signed-minor tightness conditions of Weller~\cite{Weller2016,WellerRowlandSontag2016}
read through the frame's bipartition.


\section{The structure theory}\label{sec:structure}

\subsection{The trichotomy}\label{sub:trichotomy}

Write a symbolic language as a relation $\rho$ (the legal adjacent pairs of a shift
of finite type); a protocol reports windows on the configurations legal for
$\rho$. The general problem splits into three regimes.

\begin{theorem}[acyclic reporting is tame; Vorob'ev]\label{thm:acyclic}
If the reporting hypergraph of a protocol is acyclic, the protocol is
commensurable for every language $\rho$.
\end{theorem}

This is Vorob'ev~\cite{Vorobev1962} in the present language: consistent marginals
over a regular (acyclic) complex extend, and acyclicity is the running-intersection
condition of database theory~\cite{BeeriFaginMaierYannakakis1983}, the two
identified via~\cite{Abramsky2013}. Conditional-product glue over a junction tree
assembles the window-marginals into a global distribution in the variety of
$\rho$, so frustration requires a cycle.

\begin{theorem}[full-support coordinate protocols are tame]\label{thm:product}
A $k$-context protocol of coordinate contexts with full product support is
commensurable: the product measure realises any coherent margins.
\end{theorem}

The third regime is wild: a protocol may carry a \emph{constraint variety} --- the
configurations restricted by $\rho$ --- and there the local polytope inherits
arbitrary structure.

\begin{theorem}[variety-typed protocols are universal]\label{thm:universal}
For protocols with three or more coordinate contexts subject to a constraint
variety, the achievable coherence polytopes realise every rational polytope up
to affine equivalence. Consequently: no finite facet taxonomy exists, model
membership is \textsf{NP}-complete already at three contexts, and no bounded
certificate normal form can classify contextual models in general.
\end{theorem}

\begin{proof}
The covariance map~\cite{AvisImaiItoSasaki2005} makes the polytopes
cut/correlation polytopes; by De~Loera--Onn~\cite{DeLoeraOnn2004} every rational
polytope is a face of a three-way transportation polytope, realised as imposed
zeros, i.e.\ as constraint varieties; the \textsf{NP}-completeness floor
is~\cite{Pitowsky1991}.
\end{proof}

This is why the classification is posed at the protocol level (a $\forall\exists$
question over structured families) and not the model level, where it is
intractable. The rest of the paper classifies protocol families where structure
defeats the universality.


\subsection{The parity theorem}\label{sub:parity}

The cleanest cyclic case is the golden-mean shift: $\rho_{\mathrm{gm}}$ forbids
adjacent $1$s. Observe it around a ring of length $L$ (scope contexts = the $L$
edge-windows of the cycle $C_L$).

\begin{theorem}[parity]\label{thm:parity}
The golden-mean ring protocol on $C_L$ is commensurable iff $L$ is even.
\end{theorem}

\begin{proof}
\emph{Reduction (marginal-determination on the golden-mean ring).} Write
$u_i\in[0,1]$ for the marginal $\Pr[x_i=1]$ at ring site $i\in\bbZ_L$. On the
variety $x_i=1\Rightarrow x_{i+1}=0$, so the joint distribution of the adjacent
pair $(x_i,x_{i+1})$ has cells $(0,0),(0,1),(1,0)$ with probabilities
$p_{00},p_{01},p_{10}$ and $p_{11}=0$. The endpoint marginals give
$p_{10}+p_{11}=u_i$ and $p_{01}+p_{11}=u_{i+1}$, whence, using $p_{11}=0$ and
$\sum p=1$,
\[
  p_{10}=u_i,\qquad p_{01}=u_{i+1},\qquad p_{00}=1-u_i-u_{i+1}.
\]
The edge distribution is thus \emph{affine} in $(u_i,u_{i+1})$ --- the language is
marginal-determined --- and nonnegativity of $p_{00}$ is exactly the edge
constraint $u_i+u_{i+1}\le 1$. Adjacent-context coherence is the shared marginal
$u_{i+1}$, automatic in these coordinates. Hence
\[
  C \;\cong\; \FSTAB(C_L)=\{u\in[0,1]^L : u_i+u_{i+1}\le 1\ \forall i\in\bbZ_L\},
\]
the fractional stable-set polytope of the cycle $C_L$, and
$R\cong\STAB(C_L)=\operatorname{conv}\{\mathbf 1_S : S\subseteq\bbZ_L\text{
independent}\}$, since the realisable $\{0,1\}$-marginal profiles are exactly the
indicators of independent sets (a homomorphism $C_L\to\rho_{\mathrm{gm}}$ is a
$1$-labelling with no two adjacent, i.e.\ an independent set).

\emph{Parity of the incidence determinant.} The edge constraints
$u_i+u_{i+1}\le1$ have coefficient matrix $M$, the $L\times L$ vertex--edge
incidence matrix of $C_L$ read cyclically: $M=I+P$ where $P$ is the cyclic shift
$P e_i = e_{i+1}$. Then $\det M=\det(I+P)=\prod_{k=0}^{L-1}(1+\omega^k)$ over the
$L$-th roots of unity $\omega^k$ ($P$ has eigenvalues the roots of unity), and
$\prod_{k}(1+\omega^k)=1-(-1)^L$ since $\prod_k(z-\omega^k)=z^L-1$ evaluated at
$z=-1$ gives $(-1)^L-1$, so $\prod_k(-1-\omega^k)=(-1)^L-1$ and
$\prod_k(1+\omega^k)=(-1)^L\bigl((-1)^L-1\bigr)=1-(-1)^L$. Thus
\[
  \det M \;=\; 1-(-1)^L \;=\;
  \begin{cases}0,& L\text{ even},\\ 2,& L\text{ odd}.\end{cases}
\]
Every \emph{proper} square submatrix of $M$ omits at least one ring edge, hence
is the incidence matrix of a disjoint union of \emph{paths} (a forest); a path
incidence matrix is triangularisable by ordering vertices along the path, so its
determinant lies in $\{-1,0,1\}$ --- every proper square submatrix is totally
unimodular.

\emph{Even case.} For $L$ even, $C_L$ is bipartite, so $M$ itself is totally
unimodular (bipartite incidence matrices are TU, König/Hoffman--Kruskal). With
integral right-hand sides $\FSTAB(C_L)$ has only integral vertices, and its
integral points are exactly the independent-set indicators, so
$\FSTAB(C_L)=\STAB(C_L)$: $C=R$, commensurable.

\emph{Odd case --- the unique fractional vertex.} For $L$ odd, let $u^\ast$ be a
fractional (non-integral) vertex of $\FSTAB(C_L)$. A vertex is determined by $L$
tight constraints among $\{u_i\ge0\}$ and $\{u_i+u_{i+1}\le1\}$. Suppose some
coordinate is tight at $u^\ast_j=0$. Deleting site $j$ breaks the cycle: the
remaining active constraints are those of a \emph{path} on the other sites, whose
incidence matrix is TU (forest), so every vertex of that face is integral ---
$u^\ast$ would be integral, contradiction. Hence $u^\ast>0$ coordinatewise, so no
$u_i\ge0$ constraint is tight, and all $L$ tight constraints must be the edge
constraints $u_i+u_{i+1}=1$ for every $i\in\bbZ_L$. This linear system is
$Mu=\mathbf 1$; since $L$ is odd $\det M=2\neq0$, so it has the unique solution.
Summing $u_i+u_{i+1}=1$ around the ring gives $2\sum_i u_i=L$, and the cyclic
alternation $u_{i+1}=1-u_i$ around an odd cycle forces $u_i=1-u_i$, i.e.
\[
  u^\ast \equiv \tfrac12 .
\]
Thus $\FSTAB(C_L^{\mathrm{odd}})$ is $\STAB(C_L)$ with exactly one extra vertex
$u\equiv\tfrac12$: the $\EA/\PR$ gap on an odd golden-mean ring is a single point.
\qed
\end{proof}

\begin{corollary}[the odd-ring PR-box]\label{cor:prbox}
For $L$ odd the half-model $u\equiv\tfrac12$ is separated from $R=\STAB(C_L)$ by
the odd-cycle inequality $\sum_i u_i \le (L-1)/2$, with relative violation
$1/(L-1)$; and it is strongly contextual --- its contextual fraction is $1$. It is
the golden-mean ring's analogue of a PR-box.
\end{corollary}

\begin{proof}
\emph{Separating facet.} Every independent set $S$ of the odd cycle $C_L$ has
$|S|\le(L-1)/2$ (an independent set misses at least one vertex of each of the
$\lceil L/2\rceil$ disjoint-as-possible edges; the maximum independent set of an
odd cycle has size $(L-1)/2$), so $\sum_i u_i\le(L-1)/2$ holds at every vertex of
$R$, hence on all of $R$. At $u\equiv\tfrac12$ the left side is $L/2$, exceeding
$(L-1)/2$ by $\tfrac12$; relative to the bound this is
$\tfrac12 / \tfrac{L-1}{2} = 1/(L-1)$.

\emph{Contextual fraction $1$.} The contextual fraction is $1-\lambda^\ast$, where
$\lambda^\ast=\max\{\lambda:\ u-\lambda\nu\ge0\ \text{cellwise for some }\nu\in R\}$
is the largest realisable part fitting under $u$. On each edge $u$ places mass
$\tfrac12$ on $(1,0)$, $\tfrac12$ on $(0,1)$, and $0$ on $(0,0)$. Every $\nu\in R$
is a mixture of independent-set configurations, and any configuration omitting both
endpoints of some edge puts mass on that edge's $(0,0)$ cell, which $u$ forbids. So
an admissible $\nu$ must occupy exactly one endpoint of every edge: $v_i+v_{i+1}=1$
for all $i\in\bbZ_L$. Summing gives $2\sum_i v_i=L$, impossible for integer
$\sum_i v_i$ when $L$ is odd. Hence $\lambda^\ast=0$ and $u$ is strongly
contextual. \qed
\end{proof}

\begin{remark}[dynamical reading]\label{rem:parity-dynamics}
Sliding edge-windows on a period-$L$ orbit identify time modulo $L$, so a cyclic
scope structure is recurrence and Theorem~\ref{thm:parity} says golden-mean
dynamics produces contextual data exactly on odd-period recurrence. The separating
facet is the odd-cycle inequality, one instance of the $n$-cycle noncontextuality
family~\cite{AraujoQuintinoBudroniTerraCunhaCabello2013} whose $n=5$ case is the
KCBS bound~\cite{KCBS2008}; the parity split is the golden-mean face of that
dichotomy, reached by recurrence rather than exclusivity.
\end{remark}

The spectrum varies with the language: $\{\le\}$ is tame on all rings; XOR on even
rings (odd rings have empty variety); the full language on none. More constraint
can mean more tame.


\subsection{Tamings and Circuit Localization}\label{sub:tamings}

A commensurable language keeps the winding obstruction of \S\ref{sub:winding} from
forming --- $\FSTAB=\STAB$ with no fractional vertex. It can do so in one of two
ways: by \emph{exhibiting} integrality directly (a terminal certificate), or by
\emph{stripping} structure until a terminal certificate applies (a reduction). The
mechanisms we catalogue fall under this dichotomy; each is a route to the same
one-point criterion, not an independent phenomenon.

\begin{definition}[the taming mechanisms]\label{def:tamings}
A protocol/language is tamed by one of the following. \emph{Terminal integrality
certificates:}
\begin{enumerate}[label=\textup{(\arabic*)}, itemsep=1pt]
  \item \emph{chains / empty-intersection}: the pointwise bound
    $\sum_t \mathbf 1_{c_t}\le m-1$ (\S\ref{sub:parity} base case);
  \item \emph{cross-context cliques}: perfect-graph coverage integrality;
  \item \emph{signability}: total unimodularity from a $\{\pm1\}$ column signing
    along the frame's bipartition, by Hoffman--Gale (appendix
    to~\cite{HellerTompkins1956}; cf.\ Ghouila-Houri~\cite{GhouilaHouri1962}),
    with K\"onig as the edge case;
  \item \emph{twin expansion}: a below-total-unimodularity integrality mechanism
    (a phase-forced decomposition), realised by a constructive $\theta$-sweep.
\end{enumerate}
\emph{Reductions to a terminal certificate:}
\begin{enumerate}[label=\textup{(\arabic*)}, itemsep=1pt, start=5]
  \item \emph{transient branching}: branching off the recurrent core is deleted;
  \item \emph{direct-sum decomposition}: a disconnected relation splits
    componentwise;
  \item \emph{grading}: the transfer digraph fibered over $\bbZ_d$ ---
    equivalently, Perron--Frobenius imprimitivity of period $d$;
  \item \emph{monotone collapse}: a non-symmetric relation induces monotone
    potentials that strong connectivity flattens (the SCC-reduction lemma);
  \item \emph{bipartite-target phase decoupling}: a bipartite target is a fibered
    join of two copies of its edge polytope over a global phase, reducing $C=R$ to
    the signable certificate~(3).
\end{enumerate}
The remaining case, \emph{determinism} (a functional transfer relation), is
degenerate: the polytope collapses to a point.
\end{definition}

The dichotomy is not merely presentational for two of the entries: grading is
\emph{provably} imprimitivity, and phase decoupling \emph{provably} reduces to
signability (Lemma~\ref{lem:taming10}). The others are grouped by the $C=R$
mechanism they exhibit or invoke, not by their realisation --- three of them share
the $\theta$-sweep as a construction while resting on different integrality
grounds. The signable certificate is representative.

\begin{theorem}[signable languages are tame on bipartite frames]\label{thm:taming7}
Let a language be \emph{marginal-determined} (its cell simplex is affine in the
endpoint marginals) and \emph{signable} (the edge-constraint matrix admits a
$\{\pm1\}$ column signing along the frame's bipartition). Then on every bipartite
frame it is commensurable ($C=R$). In the binary case the realisation is
moreover constructive: an anti-comonotone $\theta$-sweep (threshold $\theta$ on
one bipartition class, $1-\theta$ on the other) reproduces any coherent point as
an explicit mixture of global configurations.
\end{theorem}

\begin{proof}
There are two claims: $C=R$, and a constructive realisation.

\emph{$C=R$ (integrality).} Fix a bipartite frame with $2$-colouring
$V=V_+\sqcup V_-$. Marginal-determination lets us use vertex-marginal
coordinates; the coherence polytope $C$ is then cut out by the edge-nonnegativity
inequalities of the language together with the marginal-normalisation equalities.
Signability is the hypothesis that the coefficient matrix $M$ of the
edge-constraint rows admits a $\{\pm1\}$ column signing: multiplying the columns
indexed by $V_-$ by $-1$ turns $M$ into a matrix each of whose rows has its
nonzero entries of \emph{opposite} sign, because every edge of the frame joins a
$V_+$ vertex to a $V_-$ vertex (this is exactly where $2$-colourability enters).
A $\{0,\pm1\}$ matrix in which every column can be signed so that every row's two
nonzeros are of opposite sign is totally unimodular --- this is the sufficient
condition of Hoffman and Gale (appendix
to~\cite{HellerTompkins1956}; equivalently Ghouila-Houri's
interval-signing criterion~\cite{GhouilaHouri1962}). With integral right-hand
sides (the marginal constraints have $\{0,1\}$ data), a totally unimodular system
has only integral vertices. Each integral vertex of $C$ is a $\{0,1\}$
vertex-marginal profile, i.e.\ a global configuration, so it lies in
$R=\operatorname{conv}\{v(x)\}$; hence $C=R$.

\emph{The $\theta$-sweep realisation (binary case).} Integrality gives $C=R$
non-constructively; when the alphabet is binary the sweep exhibits, for any
coherent point $q\in C$, an explicit probability mixture of global configurations
inducing it. (For a binary marginal-determined language the edge law $q_e$ is the
countermonotone Fréchet coupling of its endpoint marginals --- e.g.\ $p_{11}=0$
with $u_i+u_{i+1}\le1$ on the golden mean --- which is exactly what the
anti-comonotone sweep produces; for larger alphabets $q_e$ need not be a monotone
coupling and the sweep is only illustrative.) Parametrise by a threshold
$\theta\in[0,1]$. On the class $V_+$ read a state off $q$'s marginals by the
\emph{comonotone} rule --- accumulate the marginal mass in a fixed state order
and select the state whose cumulative interval contains $\theta$; on $V_-$ read
by the \emph{anti-comonotone} rule with threshold $1-\theta$ (the reversed
order). For each frame edge $e=\{u_+,u_-\}$ the pair of selected endpoint states,
as $\theta$ ranges over $[0,1]$, sweeps a monotone staircase of $(V_+,V_-)$-state
pairs; because the edge distribution is affine in the endpoint marginals
(marginal-determination) and both readings are driven by the \emph{single}
scalar $\theta$, the induced pair-frequency over the uniform law
$\theta\sim\mathrm{Unif}[0,1]$ equals $q$ on that edge exactly. Coherence
(pairwise marginal agreement) guarantees the per-edge sweeps share endpoint
marginals at every shared vertex, so the $\theta$-indexed selection is a
\emph{global} configuration $x(\theta)$ for each $\theta$ at which no threshold
coincides with a jump, of which there are finitely many. The mixture
\[
  \mu \;=\; \int_0^1 \delta_{x(\theta)}\,d\theta
\]
is a finite convex combination of global configurations --- one per sub-interval
between consecutive jumps, weighted by its length --- inducing $q$. The
construction is a finite case analysis on the ordering of the jump-thresholds
followed by Lebesgue length on $[0,1]$. \qed
\end{proof}

A counting fact then closes the loopless symmetric case entirely.

\begin{theorem}[forest counting]\label{thm:forest}
A symmetric loopless language is marginal-determined iff, as a homomorphism
target, it is a forest. Every such language phase-decouples on bipartite frames.
\end{theorem}

\begin{proof}
Write $H$ for the target graph on vertex set $A$ (the states) whose edge set is
the loopless symmetric relation $\rho$; symmetry and looplessness mean
$\rho=\{(a,b):\{a,b\}\in E(H)\}$, so $|\rho|=2|E(H)|$.

\emph{Marginal-determination $\Rightarrow$ forest.} Marginal-determination asks
that the endpoint-marginal map --- edge distribution $\mapsto$ pair of endpoint
vertex-marginals --- be injective. Cells are \emph{ordered} pairs, so the map is
the directed incidence structure. A cycle in $H$ obstructs injectivity: the
directed circulation around it ($+\varepsilon$ on the forward cell, $-\varepsilon$
on the reverse cell of each cyclic edge) has zero endpoint-marginal image --- at
every vertex the forward and reverse contributions cancel in both the tail and the
head marginal --- so it is a kernel vector and a cyclic component makes the map
non-injective. (This uses looplessness and symmetry: the reverse cell of each edge
is legal.) Hence $H$ is acyclic, a forest.

\emph{Forest $\Rightarrow$ marginal-determined and bipartite.} A forest is
bipartite. On a tree the endpoint-marginal map is injective: orienting from a root,
each edge distribution is recovered as the coupling of its two endpoint marginals,
one edge at a time, with no cycle to obstruct. Bipartiteness makes the target
phase-decoupling --- a $2$-colouring of $H$ puts one endpoint of each edge in each
class --- so Theorem~\ref{thm:taming7} and phase decoupling (taming~(9)) apply. \qed
\end{proof}

The assembly turns on phase decoupling, taming~(9).

\begin{lemma}[bipartite-target phase decoupling; taming (9)]\label{lem:taming10}
Let $\rho$ be a marginal-determined language whose target graph $H$ is bipartite,
with $2$-colouring $A=A_0\sqcup A_1$. Then on every bipartite frame the coherence
polytope fibers over a single \emph{phase} scalar as
\[
  C \;=\; \bigcup_{t\in[0,1]} \bigl(t\cdot\FSTAB \;\oplus\; (1-t)\cdot\FSTAB\bigr),
  \qquad
  R \;=\; \bigcup_{t\in[0,1]} \bigl(t\cdot\STAB \;\oplus\; (1-t)\cdot\STAB\bigr),
\]
two copies of the edge polytope braided over $t$; and $\FSTAB=\STAB$ by
Theorem~\ref{thm:taming7}, so $C=R$.
\end{lemma}

\begin{proof}
A connected bipartite frame admits exactly two global colourings by $H$'s classes,
carrying masses $t$ and $1-t$; the phase $t$ is the single global degree of
freedom. Conditioned on either colouring, the residual data is an edge distribution
over $H$ in vertex-marginal coordinates, a point of $\FSTAB$, and the two are
independent given $t$ --- the fibered direct sum. Bipartiteness makes the
edge-constraint matrix signable along the frame's bipartition, so $\FSTAB=\STAB$ by
Theorem~\ref{thm:taming7} and $C=R$ termwise in $t$. \qed
\end{proof}

\begin{theorem}[Circuit Localization, symmetric loopless case]\label{thm:cl}
Every symmetric loopless marginal-determined language is commensurable on every
gate-passing frame. Contextuality is excluded by structure: such a language is a
forest target (Theorem~\ref{thm:forest}), tamed by phase-decoupling on bipartite
frames and by acyclicity (Theorem~\ref{thm:acyclic}) otherwise.
\end{theorem}

\begin{proof}
By Theorem~\ref{thm:forest} a symmetric loopless marginal-determined language is,
as a homomorphism target, a forest, hence bipartite. Let a gate-passing frame be
given. If the reporting hypergraph is acyclic, the protocol is commensurable for
\emph{every} language by Theorem~\ref{thm:acyclic}, $\rho$ included. Otherwise the
frame carries a cycle; restricted to the bipartite target, the coherence polytope
phase-decouples by Lemma~\ref{lem:taming10}, which reduces $C=R$ to the signable
case (Theorem~\ref{thm:taming7}). Every gate-passing frame falls under one of
these two cases, so the language is commensurable on all of them. The fork ---
a marginal-determined symmetric language contextual on some gate-passing frame ---
is therefore excluded by structure, not by search. \qed
\end{proof}

\begin{remark}[scope of the theorem]\label{rem:cl-boundary}
The theorem is proved for loopless targets. Two extensions are verified only to
finite scale, not as theorems: the \emph{loopy} symmetric case (no
pair-bipartite non-signable survivor appears at $\le4$ states, but larger
alphabets are open) and the exclusion of \emph{cyclic} symmetric targets
(established, over all lengths, only for golden-mean odd rings,
Theorem~\ref{thm:parity}). These two steps are owed.
\end{remark}


\subsection{The winding characterisation and its certified discriminant}\label{sub:winding}

Lift a length-$L$ ring protocol to its \emph{layered graph}: vertices $(i,a)$ for
layer $i\in\bbZ_L$ and state $a$, arcs $(i,a)\to(i{+}1,b)$ for legal pairs. The
\emph{winding} of a closed walk is its net layer-advance divided by $L$.

\begin{theorem}[winding characterisation]\label{thm:winding}
Under the genericity gate (non-adjacent contexts share only trivial events), a
length $L$ is unsafe for $\rho$ iff the layered ring over $\bbZ_L$ carries a
layer-injective simple cycle of winding $\ge 2$. The characterisation is
structural; only its eventual-periodicity refinement (\S\ref{sec:open}) is
conjectural.
\end{theorem}

The proof has three links: a gate lemma identifying $C$ with a circulation
polytope; flow decomposition reading its vertices as simple cycles; and the
layer-injectivity discriminant pinning simplicity.

\begin{lemma}[gate lemma]\label{lem:gate}
Under the genericity gate, the coherence polytope $C$ of the length-$L$ ring
protocol equals the normalised circulation polytope of the layered graph: the EA
(coherence) constraints and the flow-conservation constraints have the same
solution set. Consequently a coherent point is a normalised circulation, and
conversely.
\end{lemma}

\begin{proof}
Give each arc $(i,a)\to(i{+}1,b)$ the probability $q_i(a,b)$ of the ordered pair
at the adjacent contexts $(i,i{+}1)$. Coherence is pairwise marginal agreement:
$\sum_b q_{i-1}(b,a)=\sum_b q_i(a,b)$ for every $(i,a)$, which is Kirchhoff
conservation at vertex $(i,a)$. So a coherent point is a circulation, normalised by
unit mass per context. Conversely, a circulation is coherent iff it agrees with the
protocol on shared events; under the genericity gate the only shared events are the
adjacent overlaps, whose agreement is the conservation already imposed, so the two
systems coincide. (When the gate fails, non-adjacent contexts impose further
identifications and $C$ is a proper face of the circulation polytope.) \qed
\end{proof}

\begin{lemma}[flow decomposition of vertices]\label{lem:flowdecomp}
A vertex of the normalised circulation polytope of the layered graph is a
normalised \emph{simple} directed cycle. It is integral (hence a realisable
global configuration) iff its winding is $1$; a fractional vertex is exactly a
normalised simple directed cycle of winding $\ge 2$.
\end{lemma}

\begin{proof}
By the flow-decomposition theorem every circulation is a nonnegative sum of
simple-cycle flows. An extreme point admits no such sum with two terms, so a vertex
is a single normalised simple cycle $\gamma$. If $w(\gamma)=1$ then $\gamma$ visits
each layer once and is the graph of a configuration $x:\bbZ_L\to A$, an integral
point of $R$. If $w(\gamma)\ge2$ it visits some layer twice, spreading that layer's
mass fractionally: non-integral. So the fractional vertices are exactly the
winding-$\ge2$ simple cycles, and $C\neq R$ iff one exists. \qed
\end{proof}

A winding-$\ge2$ closed walk may revisit a vertex and fail to be simple; the
certificate of simplicity is layer-injectivity, each layer visited in distinct
states.

\begin{lemma}[layer-injectivity discriminant; Lean-certified]\label{lem:injectivity}
A layered-ring closed walk is a simple cycle (its vertex map is injective) iff it
is layer-injective (its state map is injective on each layer fibre).
\end{lemma}

The discriminant and the winding-$1$/configuration correspondence at the heart of
Lemma~\ref{lem:flowdecomp} are machine-verified in Lean~4, zero \texttt{sorry}s,
standard axioms only.\footnote{\texttt{WindingInjectivity.lean}
(\texttt{simple\_iff\_layerInjective} and consequences) and
\texttt{WindingDichotomy.lean} (\texttt{winding\_one\_isConfig},
\texttt{winding\_ge\_two\_fractional}), in
\texttt{formalization/QuerySystem/}.} What Lean certifies is the combinatorial
kernel --- a winding-$1$ trace is a global configuration, a simple trace longer
than $L$ is fractional --- \emph{given} the classical flow-decomposition theorem
(Ahuja--Magnanti--Orlin), which is a cited axiom, not proved. It does not certify
the winding characterisation itself: the polytope-vertex step of
Lemma~\ref{lem:flowdecomp} and the gate identification of Lemma~\ref{lem:gate}
remain on paper.

\begin{proof}[Proof of Theorem~\ref{thm:winding}]
By Lemma~\ref{lem:gate} $C$ is the circulation polytope; by
Lemma~\ref{lem:flowdecomp} $L$ is unsafe iff $C$ has a fractional vertex, iff a
simple cycle of winding $\ge2$ exists; by Lemma~\ref{lem:injectivity} such a cycle
is exactly a layer-injective one. \qed
\end{proof}


\section{The open case and consequences}\label{sec:open}

The symmetric case is settled (\S\ref{sub:tamings}); any remaining contextual
protocol is strongly connected and non-symmetric, the SCC-reduction lemma
(monotone collapse, taming~(8)) reducing commensurability to that class. Exhaustive search finds no contextual
protocol there at small scale (none at $\le4$ states, none through $9$ edges at
$5$).

\begin{conjecture}[universal impossibility]\label{conj:universal}
The fork does not exist: Circuit Localization holds universally for total
languages on gate-passing, time-realisable frames. Equivalently, blocking the
frustrating exchange at a rich set of lengths forces one of the catalogued
tamings, each of which localises the polytope.
\end{conjecture}

The intended route runs through eventual periodicity of the unsafe set (a
Wielandt-bounded union of arithmetic progressions), which reduces to a
combinatorial pruning lemma.

\begin{conjecture}[pruning lemma]\label{conj:pruning}
Let $\TR_k$ be the lengths admitting a diagonal-avoiding tuple-rotation walk in
the $k$-fold tensor power, and $\LISC_k$ the lengths admitting a layer-injective
simple cycle of winding $k$. Then $\TR_k=\LISC_k$ cofinitely.
\end{conjecture}

The inclusion $\LISC_k\subseteq\TR_k$ is immediate. The reverse, at $k=2$, reduces
to whether two exchange-linked winding-$1$ strands form a single simple cycle ---
but the reduction drops a relative-phase degree of freedom (a winding-$2$ simple
cycle interleaves its wraps at a phase a lockstep exchange cannot express), so the
$k=2$ residual is phase-parametrised strand-disjointness. No $\TR_k\setminus\LISC_k$
disagreement appears to layered length $35$: evidence, not proof.


\subsection{Consequences for reconstruction}\label{sub:consequences}

On the tame classes --- acyclic reporting (Theorem~\ref{thm:acyclic}), symmetric
marginal-determined languages (Theorem~\ref{thm:cl}), even-cyclic golden-mean
recurrence (Theorem~\ref{thm:parity}) --- coherent window-statistics determine a
global process, and Theorem~\ref{thm:winding} makes the boundary decidable: a
layer-injective winding-$\ge2$ cycle separates the contextual protocols from the
safe ones. This is the licensing layer beneath process reconstruction
(``statistical Takens'': a measure-plus-shift, not a manifold). By
Theorem~\ref{thm:universal} the boundary is intrinsic --- general variety-typed
protocols are as hard as rational polytope membership --- so contextuality is a
feature of observing through windows, present exactly when the geometry admits the
winding obstruction; where Conjecture~\ref{conj:universal} holds, no total language
manufactures it unless the exchange criterion fires.

The same failure of local coherence to globalise appears in Paper~II algebraically
(dispersion-free states collapsing $\PR$ on non-distributive logics) and in the
$\sigma$-essential witness set-theoretically (a finitely-coherent pattern with no
$\sigma$-additive global state). Dynamically it is sharpest: frustration is an
unwitnessed exchange --- two histories trading places around the clock of
observation without meeting --- and contextual data is the shadow of an unresolved
covering of time.


\section*{Note on method}\label{sec:method}

Each commensurability verdict is a theorem or an exact rational enumeration, each
contextuality verdict an exhibited rational witness in $C\setminus R$; sampling
serves only as screening. Lemma~\ref{lem:injectivity} and the flow-decomposition
kernel are machine-checked in Lean~4; the constructive tamings ($\theta$-sweep,
phase-mixture, monotone collapse) are the remaining certification targets.
